%==============================================================
% Notater til Ch3 Methods — mesh / y+
%==============================================================

% - Nevn at TurboGrid brukes på blader (RDT, GV) og unstructured + inflation
%   på rør/draft tube
% - SST k-omega med automatic wall treatment over hele domenet
% - Konsistent mesh-strategi på tvers av alle 18 case
% - Rapporter y+ statistikk i tabell (gjennomsnitt over de 18 case):
%     RDT-blader: mean y+ ~15-25, max ~30
%     GV-blader:  mean y+ ~20-25, max ~70
%     DT-vegg:    mean y+ ~170
% - Forklar at blade-mesh ligger i buffer-sonen (5 < y+ < 30),
%   ikke i resolved viscous sublayer
% - Begrunn valget: y+ <= 1 ville krevd 2-3x flere celler per case
%   --> ikke praktisk for 18-case screening
% - Pek på konsistens med tidligere oppgaver på samme rig
%   (Dokset, Andersen, Jensen) — alle i samme y+ range
% - Konkluder: integral metrics (dH, swirl, ranking) er pålitelige
%   til tross for buffer-sone


%==============================================================
% Notater til Ch6 / Future Work — mesh refinement
%==============================================================

% - Anbefal refinement til y+ <= 1 på blader for cavitation-fokus
% - Begrunn hvorfor det er viktig: lokale trykk-minima på RDT-suge-side
%   krever resolved viscous sublayer
% - Praktisk anbefaling: gjør det kun for de mest lovende case
%   fra screening, ikke alle 18
% - Nevn at det er kostbart — ~2-3x cellecount per case
% - Mulig kobling til transient/URANS hvis mesh refines uansett


%==============================================================
% Stikkord til selve skrivingen
%==============================================================

% - Hold det kort — y+ er ikke hovedtema
% - Én tabell, ett avsnitt i Methods, én setning i Future Work holder
% - Ikke unnskyld meshet — begrunn det med kost/nytte for 18-case screening
% - Vis at du har sjekket det kvantitativt (tabellen viser det)